\begin{afterword}
\summaryhead

The post defends narrow words and categories. In their own trades people draw fine distinctions: a mechanic does not lump a carburetor and a radiator together as ``car parts.'' Outside their trades people often stretch words as wide as they will go, because general theories sound loftier.

A parable makes the point. The curious pick out a few pebbles and call them diamonds; others, meaning well, want the name for every pebble. Even poets must draw fine distinctions. Biologists may explain living things without explaining the ``evolution'' of stars or technology. ``Everything is connected to everything else'' says nothing, because a graph with every edge carries no more information than a graph with none. Good categories and good hypotheses exclude something, and Newton was right to explain gravity and not money. Sneering at narrowness recalls Plato's contempt for learning by looking. The post ends by insisting that there is nothing wrong with narrow categories.

\responsehead

The post's central point is sound. A category that includes everything, like a hypothesis that permits every outcome, tells you nothing, and the argument that a graph with every possible edge carries no more information than a graph with none makes this cleanly. For hypotheses the point is Popper's, as the notes on ``Your Strength as a Rationalist'' discuss.

What the post does not address is the question a reader would need answered: how to tell a good broad category from a bad one. Its parable settles that question by stipulation. The pebbles the curious pick really are diamonds, so the people who want to widen the word are wrong by construction. In any real dispute over a word, whether the other pebbles share the inner qualities is exactly what is in doubt.

The examples do not supply the missing test. The post names no one who reasons from the shared word ``evolution'' to the conclusion that ``it must all be the same thing,'' and the shared word proves nothing by itself: ``stellar evolution'' is the astronomers' own term, and the word's general sense of development is older than the biological one. The post's best example works against it. Newton's achievement was a connection, between the fall of bodies on Earth and the motion of the planets and tides. Two paragraphs earlier the post had said that detailed study shows how things are ``not alike'' and leads to ``subtracting edges.''

One smaller point. The passage from Plato is about astronomy done by gazing instead of by mathematics; it says nothing about manual labor or slaves, the link the post draws.

In short: a sound point about categories that exclude nothing, argued by a parable that settles its question by stipulation, against opponents the post does not name, with no test for which connections to keep, and with Newton as its best example, whose achievement was a connection.
\end{afterword}
